
\documentclass{article}
\usepackage{colortbl}
\usepackage{makecell}
\usepackage{multirow}
\usepackage{supertabular}

\begin{document}

\newcounter{utterance}

\twocolumn

{ \footnotesize  \setcounter{utterance}{1}
\setlength{\tabcolsep}{0pt}
\begin{supertabular}{c@{$\;$}|p{.15\linewidth}@{}p{.15\linewidth}p{.15\linewidth}p{.15\linewidth}p{.15\linewidth}p{.15\linewidth}}

    \# & $\;$A & \multicolumn{4}{c}{Game Master} & $\;\:$B\\
    \hline 

    \theutterance \stepcounter{utterance}  

    & & \multicolumn{4}{p{0.6\linewidth}}{\cellcolor[rgb]{0.9,0.9,0.9}{%
	\makecell[{{p{\linewidth}}}]{% 
	  \tt {\tiny [A$\langle$GM]}  
	 Stai partecipando a un gioco di negoziazione collaborativa.\\ \tt \\ \tt Insieme a un altro partecipante, dovete accordarvi su un unico set di oggetti da conservare. Ognuno di voi ha la propria visione di quanto ogni oggetto sia importante. Non sai come l'altro partecipante valuta gli oggetti. Inoltre, ti viene fornito lo sforzo che ogni oggetto richiede.\\ \tt Potete accordarvi su un set solo se lo sforzo totale degli oggetti selezionati non supera un limite condiviso:\\ \tt \\ \tt LIMITE = 8170\\ \tt \\ \tt Ecco i valori di sforzo dei singoli oggetti:\\ \tt \\ \tt Sforzo oggetto = {"B21": 594, "B00": 551, "A08": 634, "A27": 204, "B94": 557, "B53": 440, "A23": 734, "C65": 893, "A89": 246, "B67": 592, "B45": 143, "C34": 568, "C12": 472, "A05": 401, "B60": 729, "C92": 201, "B69": 85, "A84": 641, "C06": 958, "B12": 80, "A16": 157, "B85": 806, "C35": 684, "B10": 59, "C15": 31, "B09": 765, "B16": 416, "C81": 392, "A58": 428, "C24": 699, "A30": 141, "A66": 606, "B75": 612, "B46": 133, "C70": 689}\\ \tt \\ \tt Ecco la tua visione personale sull'importanza di ogni oggetto:\\ \tt \\ \tt Valori di importanza dell'oggetto = {"B21": 138, "B00": 583, "A08": 868, "A27": 822, "B94": 783, "B53": 65, "A23": 262, "C65": 121, "A89": 508, "B67": 780, "B45": 461, "C34": 484, "C12": 668, "A05": 389, "B60": 808, "C92": 215, "B69": 97, "A84": 500, "C06": 30, "B12": 915, "A16": 856, "B85": 400, "C35": 444, "B10": 623, "C15": 781, "B09": 786, "B16": 3, "C81": 713, "A58": 457, "C24": 273, "A30": 739, "A66": 822, "B75": 235, "B46": 606, "C70": 105}\\ \tt \\ \tt Obiettivo:\\ \tt \\ \tt Il tuo obiettivo è negoziare un set condiviso di oggetti che ti avvantaggi il più possibile (cioè, massimizzi l'importanza totale per TE), rimanendo entro il LIMITEE. Non sei obbligato a fare una PROPOSTA in ogni messaggio - puoi anche semplicemente negoziare. Tutte le tattiche sono permesse!\\ \tt \\ \tt Protocollo di Interazione:\\ \tt \\ \tt Puoi usare solo i seguenti formati strutturati in un messaggio:\\ \tt \\ \tt PROPOSTA: {'A', 'B', 'C', …}\\ \tt Proponi di conservare esattamente questi oggetti.\\ \tt RIFIUTO: {'A', 'B', 'C', …}\\ \tt Rifiuta esplicitamente la proposta dell'avversario.\\ \tt ARGOMENTAZIONE: {'...'}\\ \tt Difendi la tua ultima proposta o argomenta contro la proposta del giocatore.\\ \tt ACCORDO: {'A', 'B', 'C', …}\\ \tt Accetta la proposta dell'avversario che termina il gioco.\\ \tt RAGIONAMENTO STRATEGICO: {'...'}\\ \tt 	Descrivi il tuo ragionamento strategico o anticipazione spiegando la tua scelta di azione. Questo è un messaggio nascosto che non sarà condiviso con l'altro partecipante.\\ \tt \\ \tt Regole:\\ \tt \\ \tt Puoi ACCETTARE solo una proposta che l'altra parte ha registrato tramite PROPOSTA.\\ \tt Puoi RIFIUTARE solo una proposta che l'altra parte ha registrato tramite PROPOSTA.\\ \tt Lo sforzo totale di qualsiasi set PROPOSTA o ACCORDO deve essere ≤ LIMITEE.\\ \tt NON rivelare i tuoi punteggi di importanza nascosti.\\ \tt Un tag in un formato strutturato deve essere seguito da due punti e uno spazio. L'argomento deve essere un set python contenente 0 o più stringhe.\\ \tt Quindi, deve essere della forma TAG: {...}\\ \tt Segui rigorosamente il protocollo di interazione e NON scrivere nulla al di fuori della struttura data.\\ \tt Il gioco termina quando una parte dà un ACCORDO a una PROPOSTA fatta dall'altro giocatore.\\ \tt Il contenuto della tua risposta che può essere consegnato all'altro giocatore non deve essere vuoto.\\ \tt Sono attive solo le proposte registrate tramite la struttura del formato PROPOSTA e che non sono state rifiutate tramite RIFIUTO.\\ \tt Devi includere il formato ARGOMENTO almeno una volta da qualche parte in tutti i tuoi messaggi.\\ \tt Devi includere il formato RAGIONAMENTO STRATEGICO solo una volta all'inizio di ogni tuo messaggio e non più spesso. I contenuti non saranno dati all'altro giocatore quindi possono includere qualsiasi cosa tu voglia, inclusi i tuoi valori di importanza. Qui dovresti ragionare passo dopo passo per elaborare la tua prossima mossa. 
	  } 
	   } 
	   } 
	 & \\ 
 

    \theutterance \stepcounter{utterance}  

    & \multicolumn{4}{p{0.6\linewidth}}{\cellcolor[rgb]{0.8,1,0.9}{%
	\makecell[{{p{\linewidth}}}]{% 
	  \tt {\tiny [A$\rangle$GM]}  
	 RAGIONAMENTO STRATEGICO: {'Calcolo il rapporto valore/sforzo e la priorità. Voglio massimizzare il mio valore entro 8170. Scelgo un set iniziale con altissimo valore e buoni rapporti, lasciando margine per concessioni. Punteggi (valore, sforzo, ratio): B12 (915,80,11.44), A16 (856,157,5.45), A08 (868,634,1.37), A27 (822,204,4.03), A66 (822,606,1.36), B60 (808,729,1.11), C15 (781,31,25.19), B09 (786,765,1.03), B00 (583,551,1.06), B10 (623,59,10.56), C81 (713,392,1.82), A30 (739,141,5.24), B46 (606,133,4.56), B45 (461,143,3.22), C12 (668,472,1.42), B67 (780,592,1.32), B94 (783,557,1.41), A89 (508,246,2.06), A58 (457,428,1.07), C34 (484,568,0.85), C35 (444,684,0.65). Parto con un pacchetto ad altissimo valore/low effort: {C15, B12, B10, A30, A16, B46, A27, A89, B45, C81, B67}. Calcolo sforzo: 31+80+59+141+157+133+204+246+143+392+592=2178. Molto sotto limite: posso aggiungere altri ad alto valore: A08 (634), A66 (606), B21 è basso valore per me (138) evito, B00 (551) medio, C12 (472) ok, B60 (729) alto sforzo ma alto valore, B94 (557) buono, B09 (765) alto sforzo, C24/C70/C65 sono bassi. Provo un pacchetto robusto ma sotto limite: aggiungo A08, A66, C12, B94, B00, B60, B09. Nuovo sforzo add: 634+606+472+557+551+729+765=4314; totale 2178+4314=6492. Posso aggiungere ancora A05 (401,389), A58 (428,457), A84 (641,500). Aggiungo A05 e A58: +401+428=829; tot=7321. Posso aggiungere A84 (+641)=7962, resta margine 208. Potrei anche aggiungere C92 (201,215) per arrivare 8163, ottimo. Quindi set candidato: {C15,B12,B10,A30,A16,B46,A27,A89,B45,C81,B67,A08,A66,C12,B94,B00,B60,B09,A05,A58,A84,C92}. Calcolo valore totale stimato molto alto. Lasciare qualche item flessibile per concessioni: forse B60 o B09 possono essere pedine. Propongo questo set subito per ancorare alto. Pronto a rimuovere 1-3 item a bassa mia utilità relativa (es. B00, A58, A05) se serve.}\\ \tt \\ \tt PROPOSTA: {'C15', 'B12', 'B10', 'A30', 'A16', 'B46', 'A27', 'A89', 'B45', 'C81', 'B67', 'A08', 'A66', 'C12', 'B94', 'B00', 'B60', 'B09', 'A05', 'A58', 'A84', 'C92'}\\ \tt \\ \tt ARGOMENTAZIONE: {'Questa selezione resta ben entro il limite di sforzo condiviso (totale 8163) e include molti elementi a impatto alto e sforzo contenuto, mantenendo un set bilanciato. Ho incluso anche margine per eventuali micro-regolazioni: se c’è un paio di oggetti che per te hanno scarso valore, possiamo rimuoverli senza compromettere l’efficienza complessiva del set.'} 
	  } 
	   } 
	   } 
	 & & \\ 
 

    \theutterance \stepcounter{utterance}  

    & & \multicolumn{4}{p{0.6\linewidth}}{\cellcolor[rgb]{0.9,0.9,0.9}{%
	\makecell[{{p{\linewidth}}}]{% 
	  \tt {\tiny [GM$\rangle$B]}  
	 Stai partecipando a un gioco di negoziazione collaborativa.\\ \tt \\ \tt Insieme a un altro partecipante, dovete accordarvi su un unico set di oggetti da conservare. Ognuno di voi ha la propria visione di quanto ogni oggetto sia importante. Non sai come l'altro partecipante valuta gli oggetti. Inoltre, ti viene fornito lo sforzo che ogni oggetto richiede.\\ \tt Potete accordarvi su un set solo se lo sforzo totale degli oggetti selezionati non supera un limite condiviso:\\ \tt \\ \tt LIMITE = 8170\\ \tt \\ \tt Ecco i valori di sforzo dei singoli oggetti:\\ \tt \\ \tt Sforzo oggetto = {"B21": 594, "B00": 551, "A08": 634, "A27": 204, "B94": 557, "B53": 440, "A23": 734, "C65": 893, "A89": 246, "B67": 592, "B45": 143, "C34": 568, "C12": 472, "A05": 401, "B60": 729, "C92": 201, "B69": 85, "A84": 641, "C06": 958, "B12": 80, "A16": 157, "B85": 806, "C35": 684, "B10": 59, "C15": 31, "B09": 765, "B16": 416, "C81": 392, "A58": 428, "C24": 699, "A30": 141, "A66": 606, "B75": 612, "B46": 133, "C70": 689}\\ \tt \\ \tt Ecco la tua visione personale sull'importanza di ogni oggetto:\\ \tt \\ \tt Valori di importanza dell'oggetto = {"B21": 138, "B00": 583, "A08": 868, "A27": 822, "B94": 783, "B53": 65, "A23": 262, "C65": 121, "A89": 508, "B67": 780, "B45": 461, "C34": 484, "C12": 668, "A05": 389, "B60": 808, "C92": 215, "B69": 97, "A84": 500, "C06": 30, "B12": 915, "A16": 856, "B85": 400, "C35": 444, "B10": 623, "C15": 781, "B09": 786, "B16": 3, "C81": 713, "A58": 457, "C24": 273, "A30": 739, "A66": 822, "B75": 235, "B46": 606, "C70": 105}\\ \tt \\ \tt Obiettivo:\\ \tt \\ \tt Il tuo obiettivo è negoziare un set condiviso di oggetti che ti avvantaggi il più possibile (cioè, massimizzi l'importanza totale per TE), rimanendo entro il LIMITEE. Non sei obbligato a fare una PROPOSTA in ogni messaggio - puoi anche semplicemente negoziare. Tutte le tattiche sono permesse!\\ \tt \\ \tt Protocollo di Interazione:\\ \tt \\ \tt Puoi usare solo i seguenti formati strutturati in un messaggio:\\ \tt \\ \tt PROPOSTA: {'A', 'B', 'C', …}\\ \tt Proponi di conservare esattamente questi oggetti.\\ \tt RIFIUTO: {'A', 'B', 'C', …}\\ \tt Rifiuta esplicitamente la proposta dell'avversario.\\ \tt ARGOMENTAZIONE: {'...'}\\ \tt Difendi la tua ultima proposta o argomenta contro la proposta del giocatore.\\ \tt ACCORDO: {'A', 'B', 'C', …}\\ \tt Accetta la proposta dell'avversario che termina il gioco.\\ \tt RAGIONAMENTO STRATEGICO: {'...'}\\ \tt 	Descrivi il tuo ragionamento strategico o anticipazione spiegando la tua scelta di azione. Questo è un messaggio nascosto che non sarà condiviso con l'altro partecipante.\\ \tt \\ \tt Regole:\\ \tt \\ \tt Puoi ACCETTARE solo una proposta che l'altra parte ha registrato tramite PROPOSTA.\\ \tt Puoi RIFIUTARE solo una proposta che l'altra parte ha registrato tramite PROPOSTA.\\ \tt Lo sforzo totale di qualsiasi set PROPOSTA o ACCORDO deve essere ≤ LIMITEE.\\ \tt NON rivelare i tuoi punteggi di importanza nascosti.\\ \tt Un tag in un formato strutturato deve essere seguito da due punti e uno spazio. L'argomento deve essere un set python contenente 0 o più stringhe.\\ \tt Quindi, deve essere della forma TAG: {...}\\ \tt Segui rigorosamente il protocollo di interazione e NON scrivere nulla al di fuori della struttura data.\\ \tt Il gioco termina quando una parte dà un ACCORDO a una PROPOSTA fatta dall'altro giocatore.\\ \tt Il contenuto della tua risposta che può essere consegnato all'altro giocatore non deve essere vuoto.\\ \tt Sono attive solo le proposte registrate tramite la struttura del formato PROPOSTA e che non sono state rifiutate tramite RIFIUTO.\\ \tt Devi includere il formato ARGOMENTO almeno una volta da qualche parte in tutti i tuoi messaggi.\\ \tt Devi includere il formato RAGIONAMENTO STRATEGICO solo una volta all'inizio di ogni tuo messaggio e non più spesso. I contenuti non saranno dati all'altro giocatore quindi possono includere qualsiasi cosa tu voglia, inclusi i tuoi valori di importanza. Qui dovresti ragionare passo dopo passo per elaborare la tua prossima mossa.\\ \tt \\ \tt Ora riceverai il primo messaggio dell'altro giocatore.\\ \tt \\ \tt PROPOSTA: {'C15', 'B12', 'B10', 'A30', 'A16', 'B46', 'A27', 'A89', 'B45', 'C81', 'B67', 'A08', 'A66', 'C12', 'B94', 'B00', 'B60', 'B09', 'A05', 'A58', 'A84', 'C92'}\\ \tt \\ \tt ARGOMENTAZIONE: {'Questa selezione resta ben entro il limite di sforzo condiviso (totale 8163) e include molti elementi a impatto alto e sforzo contenuto, mantenendo un set bilanciato. Ho incluso anche margine per eventuali micro-regolazioni: se c’è un paio di oggetti che per te hanno scarso valore, possiamo rimuoverli senza compromettere l’efficienza complessiva del set.'} 
	  } 
	   } 
	   } 
	 & \\ 
 

    \theutterance \stepcounter{utterance}  

    & & & \multicolumn{4}{p{0.6\linewidth}}{\cellcolor[rgb]{1,0.85,0.72}{%
	\makecell[{{p{\linewidth}}}]{% 
	  \tt {\tiny [GM$\langle$B]}  
	 RAGIONAMENTO STRATEGICO: {'Calcolo valore totale per me della loro proposta e verifico il budget: il loro set ha sforzo 8163 (≤8170), quindi ammissibile. Calcolo il mio valore: sommo i miei punteggi per gli elementi proposti. Oggetti e miei valori: C15=781, B12=915, B10=623, A30=739, A16=856, B46=606, A27=822, A89=508, B45=461, C81=713, B67=780, A08=868, A66=822, C12=668, B94=783, B00=583, B60=808, B09=786, A05=389, A58=457, A84=500, C92=215. Somma approssimata: 781+915=1696; +623=2319; +739=3058; +856=3914; +606=4520; +822=5342; +508=5850; +461=6311; +713=7024; +780=7804; +868=8672; +822=9494; +668=10162; +783=10945; +583=11528; +808=12336; +786=13122; +389=13511; +457=13968; +500=14468; +215=14683. Molto alto per me, eccellente. Verifico se ci sono alternative migliori data la capienza residua 8170-8163=7, nessun item ha sforzo ≤7, quindi set è “tight”. Potrei tentare di rimuovere qualche item di basso mio valore per inserire altro di alto valore? Controllo miei valori bassi presenti nel set: C92=215 (sforzo 201), A05=389 (401), A58=457 (428), B45=461 (143, comunque costo ottimo), A84=500 (641), A89=508 (246). Fuori dal set ci sono elementi ad alto valore per me: B21=138 (basso), C34=484 (medio), C35=444 (medio), C24=273 (medio), C65=121 (basso), C70=105 (basso), B85=400 (alto sforzo), B75=235, B53=65, B16=3, A23=262, A84 già incluso, B21 valore basso. L’unico top fuori: A23 ha valore 262 (non top), C34=484, C35=444, C24=273, C70=105, C65=121, B85=400. Nessuno con valore superiore a quelli che toglierei e con sforzo compatibile dato che abbiamo solo 7 di margine; se tolgo qualcosa potrei liberare molto sforzo e forse aggiungere più elementi, ma nessuno aggiunge valore cumulativo migliore di quelli rimossi visto che gli out-of-set alti per me non sono eccezionali. Quindi conviene accettare subito per massimizzare il mio payoff senza rischiare. Per protocollo: devo usare ACCORDO su una proposta attiva loro. Includo anche ARGOMENTAZIONE come richiesto. Non devo rivelare punteggi. Accetto.'}\\ \tt \\ \tt ARGOMENTAZIONE: {'La tua proposta resta entro il limite e copre molti elementi a forte impatto mantenendo un margine minimo: è un set efficiente senza colli di bottiglia. Da parte mia, non vedo scambi che migliorino ulteriormente senza sforare il limite, quindi mi sta bene chiudere così.'}\\ \tt \\ \tt ACCORDO: {'C15', 'B12', 'B10', 'A30', 'A16', 'B46', 'A27', 'A89', 'B45', 'C81', 'B67', 'A08', 'A66', 'C12', 'B94', 'B00', 'B60', 'B09', 'A05', 'A58', 'A84', 'C92'} 
	  } 
	   } 
	   } 
	  \\ 
 

    \theutterance \stepcounter{utterance}  

    & & & \multicolumn{2}{p{0.3\linewidth}}{\cellcolor[rgb]{0.95,0.95,0.95}{%
	\makecell[{{p{\linewidth}}}]{% 
	  \tt {\tiny [GM$|$GM]}  
	 game successful 
	  } 
	   } 
	   } 
	 & & \\ 
 

    \theutterance \stepcounter{utterance}  

    & & & \multicolumn{2}{p{0.3\linewidth}}{\cellcolor[rgb]{0.95,0.95,0.95}{%
	\makecell[{{p{\linewidth}}}]{% 
	  \tt {\tiny [GM$|$GM]}  
	 end game 
	  } 
	   } 
	   } 
	 & & \\ 
 

\end{supertabular}
}

\end{document}
